\chapter{Introduction}
\label{chap:introduction}

\section{Nightmare Disorder and Treatment}
\label{sec:intro-irt}
% ~5% prevalence; strong comorbidity with PTSD and depression
% IRT as evidence-based gold standard -- rigid, phase-driven protocol
%   Why IRT? Its fixed structure makes it uniquely testable for protocol adherence
% Rescripting as the cognitive core -- narrative transformation as the key intervention
% Clinician shortage as the practical motivation for scalable AI-assisted delivery
Recurring nightmares are a common clinical problem with an effective
treatment: \gls{irt}, a structured protocol in which the patient rescripts
a nightmare under therapist guidance \autocite{FGA5H3ZD}. A persistent
shortage of trained therapists limits its availability
\autocite{7M7NHYXT}, making scalable, AI-assisted delivery a practical
goal.

The rescripting step, however, requires active strategy decisions
from the therapist. An AI agent that selects different strategies on
repeated runs of the same case is clinically unreliable, regardless of
how capable any single response may be. Patients also differ in how they
present: an anxious patient and a resistant patient pose different
challenges, and a model that is stable for one profile may not be stable
for another.

% Figure for IRT stages is now in Ch3 (dialogue_generation.png)

\section{AI-Assisted Psychotherapy}
\label{sec:intro-ai-therapy}
% Growing LLM deployment in mental health -- from empathy tools to full CBT state
%   machines (Stade et al., 2024; NPJ MHR, 2024)
% reDreamAI: consumer-facing IRT chatbot guiding users through all three protocol stages
%   Not a research prototype -- a real deployment target where reliability is non-negotiable
% The problem this creates: reDreamAI is not a general assistant -- it is a
%   protocol-driven agent bound to a manualized intervention
%   - A general assistant has no therapeutic obligation; variability is fine
%   - A protocol agent must make consistent strategic decisions -- inconsistency
%     is not stylistic variation, it is a clinical failure
%   - Unlike rule-based chatbots whose behaviour can be reviewed once before deployment,
%     an LLM agent must be evaluated repeatedly -- each run is a new draw from the
%     output distribution
%   - Stability is a necessary precondition for clinical deployment -- not sufficient
%     on its own, but without it, no claim about therapeutic quality is meaningful
% EU AI Act and data sovereignty requirements push toward self-hosted EU-resident
%   models -- introduces the sovereign model constraint

Large language models are increasingly deployed in mental health, from empathy-based conversational tools to protocol-driven agents that deliver structured interventions such as cognitive behavioural therapy \autocite{DTKNEPHQ,R3Y7UHDE}.
reDreamAI is the deployment context for this study: a consumer-facing IRT chatbot that guides users through all stages of the protocol.
It is not a research prototype but a real deployment target where reliability is non-negotiable.

This deployment context creates a specific problem.
A general-purpose assistant has no therapeutic obligation. Response variability is acceptable and even desirable.
A protocol-driven agent, by contrast, must make consistent strategic decisions.
If the model selects different rescripting strategies for the same patient on different runs, the patient receives inconsistent care.
Unlike rule-based chatbots whose responses can be reviewed once before deployment \autocite{YE8KQTNH}, an LLM agent generates a new response on every run.
Each run is a new draw from the output distribution, so stability must be evaluated repeatedly rather than verified once.

IRT targets nightmares rather than psychiatric diagnoses, which lowers the clinical risk relative to full psychotherapy.
Ethical considerations are discussed in \autoref{sec:disc-framework-design} and \autoref{sec:concl-limitations}.
The EU AI Act classifies autonomous therapeutic AI as high-risk and requires auditability \autocite{X5TYM56Q,A823HQRV}.
Data sovereignty requirements push toward EU-resident models \autocite{2QMKF3RM}, introducing a constraint on model selection discussed in \autoref{sec:meth-models}.

\section{The Evaluation Gap}
\label{sec:intro-eval-gap}
% Standard NLP benchmarks optimise for capability, not clinical reliability
% No existing metrics for: intent stability, protocol adherence, strategic consistency
%   The core problem: a capable but inconsistent model is clinically unsafe --
%   these two properties are not the same and are currently conflated
% Prior therapeutic AI evaluation is largely qualitative, post-hoc, or single-session
% The case for systematic, reproducible in silico clinical trials as a validation paradigm
% [FIG] Gap diagram -- what capability benchmarks measure vs. what clinical
%   reliability requires; the unmeasured space this thesis addresses

Standard \gls{nlp} benchmarks measure what a model \emph{can} do, not whether it does the same thing reliably.
A model can score well on a single evaluation while producing different therapeutic strategies on repeated runs of the same input.
Capability and stability are distinct properties, but current benchmarks conflate them.

Prior therapeutic AI evaluations, including the closest existing work
(BOLT, \cite{B97RI7GA}), measure whether a model applies appropriate
techniques. They evaluate single responses rather than the distribution
across repeated runs \autocite{JX7N8VWS}, so they cannot detect
cross-run inconsistency (\autoref{sec:bg-llm-healthcare}).
Systematic, reproducible in silico trials offer a way to measure stability at scale without requiring human participants.

% Gap argument made in prose above; no diagram needed

\section{Research Objectives}
\label{sec:intro-objectives}
% RQ: How consistent is LLM clinical reasoning across stochastic runs within
%   a structured therapeutic protocol?
% Contribution 1: A three-level hierarchical evaluation framework for therapeutic
%   AI stability -- plan consistency, response similarity, plan-response coherence
% Contribution 2: Quantitative cross-model comparison across sovereign,
%   proprietary, and open-weight model classes
% Scope: IRT rescripting phase -- isolated as the highest-stakes, most
%   cognitively demanding intervention point
%   Why isolated? Eliminates stage-routing confounds; enables clean attribution
%   of any measured instability to the therapist model

This thesis addresses the following research question:

\begin{quote}
How consistent are LLM therapeutic responses across independent runs when given the same patient context within a structured IRT protocol?
\end{quote}

\noindent Two contributions follow from this question.

\textbf{Contribution~1.}
A three-level evaluation framework that measures stability at each layer: \emph{plan consistency} (do strategy selections repeat?), \emph{response similarity} (do the generated texts convey the same meaning?), and \emph{plan-response coherence} (does the model do what it declared?).
Together, these layers reveal divergence patterns that no single metric can detect.

\textbf{Contribution~2.}
A cross-model comparison that applies this framework to sovereign, proprietary, and open-weight models under controlled temperature conditions.

The study isolates the IRT rescripting phase because it is the stage that requires active strategy decisions.
Isolating a single stage eliminates stage-routing confounds and allows any measured instability to be attributed to the therapist model rather than to upstream variation in the conversation.
\autoref{chap:methods} describes the experimental design in detail.
